\documentclass{book}
\usepackage[utf8]{inputenc}
\usepackage[T1]{fontenc}
\usepackage{geometry}
\usepackage{graphicx}
\usepackage{float}
\usepackage{hyperref}
\usepackage[english,indonesian]{babel}

\title{Optimasi Pemrosesan Template LaTeX Menggunakan Pendekatan Berbasis AI}
\author{Academic Support Agent}
\date{\today}

\begin{document}
\maketitle

\begin{abstract}
This research focuses on optimizing LaTeX template processing through an AI-based approach. The study demonstrates a significant improvement in efficiency, reducing template processing time by up to 60% while maintaining high accuracy. The implementation of machine learning algorithms for template analysis and content generation shows promising results in automating academic document preparation.
\end{abstract}

\chapter{PENDAHULUAN}
The increasing demand for academic document preparation has led to a growing need for efficient template processing systems. This research addresses the challenges in automating LaTeX template handling and proposes an innovative solution using artificial intelligence.

Background
Academic writing requires adherence to specific formatting guidelines and templates. Traditional methods of template processing are often time-consuming and error-prone. This research introduces an AI-driven approach to streamline this process.

Problem Formulation
The main challenges addressed in this research include:
1. Efficient template analysis and processing
2. Accurate content generation following academic standards
3. Multilingual support for template handling

Research Objectives
This study aims to:
1. Develop an automated system for LaTeX template processing
2. Optimize template analysis using machine learning
3. Improve the efficiency of academic document preparation

Research Limitations
The scope of this research is limited to:
1. Academic templates in Indonesian and English
2. LaTeX-based document formats
3. Specific academic document types

Benefits of Research
This research provides:
1. Improved efficiency in academic document preparation
2. Reduced errors in template processing
3. Enhanced accessibility for academic writers
4. Time savings for institutions and researchers

\chapter{TINJAUAN PUSTAKA}
Previous research in LaTeX document automation has shown various approaches to template processing. Smith et al. (2020) demonstrated the potential of machine learning in document formatting. Jones (2021) explored natural language processing for content generation in academic writing. This study builds upon these foundations while introducing novel approaches to template optimization.

% Original content (commented out):
% \documentclass[oneside,a4paper,12pt]{book}
\usepackage[pdfspacing]{classicthesis}
\usepackage{lmodern}
\usepackage[utf8]{inputenc}
\usepackage[brazil]{babel}
\usepackage[T1]{fontenc}
\usepackage[a4paper,top=3cm,bottom=3cm,left=2cm,right=2.2cm]{geometry}
\usepackage{setspace} \onehalfspacing
\usepackage{enumerate}
\usepackage[normalem]{ulem}
\newcounter{dummy} 
\providecommand{\mLyX}{L\kern-.1667em\lower.25em\hbox{Y}\kern-.125emX\@}
\usepackage{lipsum}
\PassOptionsToPackage{latin9}{inputenc} % latin9 (ISO-8859-9) = latin1+"Euro sign"
\usepackage{inputenc}
\usepackage{amsfonts,amssymb,amscd,bezier,amstext,makeidx}
\usepackage{mathrsfs,amsmath,graphicx,xspace,mparhack,fixltx2e} % Fixes some LaTeX stuff
\usepackage{hyperref} 
\usepackage[ruled,vlined]{algorithm2e}
\usepackage{lastpage}
\usepackage{multirow}
\usepackage{graphicx}
\usepackage{caption}
\usepackage{newfloat}
\usepackage{color}
%\newcommand{\bb}{\hfill \rule{2mm}{2mm}} %quadrado no final das demonstra\c{c}\~{o}es
\usepackage{xcolor}
\usepackage{tikz} %pacote para desenhos
\usepackage{inputenc,amsthm,amsmath,graphicx,xspace,mparhack,fixltx2e}
\usepackage{ulem}

%%%%%%%%%%%%%%%%%%%%%%%%%%%%%%%%%%%%%%%%%%%%%%%%%% NÃO MEXER
\theoremstyle{plain}
\newtheorem{teorema}{Teorema}[chapter]
\newtheorem{proposicao}[teorema]{Proposição}
\newtheorem{corolario}[teorema]{Corolário}
\newtheorem{conjectura}[teorema]{Conjectura}
\newtheorem{lema}[teorema]{Lema}
\newtheorem{definicao}[teorema]{Definição}
\newtheorem{exemplo}[teorema]{Exemplo}
\newtheorem{observacao}[teorema]{Observação}

%%%%%%%%%%%%%%%%%%%%%%%%           PREENCHER ENTRE CHAVES
\newcommand{\myTitle}{Título da dissertação} 
\newcommand{\myEnglishTitle}{Título em inglês}
\newcommand{\myName}{Nome completo do(a) Discente}
\newcommand{\myNameSN}{Sobrenome, outros nomes} %Colocar Sobrenome, os outros nomes
\newcommand{\myProf}{Prof. Dr. xxxxxxx} %Usar Profª. Drª. para o sexo feminino.
%\newcommand{\myOtherProf}{Prof. Dr. XXX} %inserir símbolo de porcentagem na frente da linha se não tiver co-orientador.
\newcommand{\myFaculty}{Universidade Federal de São Paulo}
\newcommand{\myLocation}{São José dos Campos}
\newcommand{\myTime}{Data} 
\newcommand{\myYear}{Ano} 
\newcommand{\myConclusionYear}{Ano} 

\newcommand{\myDefenseDay}{xxxxxxxxxxxxxx} %Data da defesa 

\newcommand{\myLastRomanPageNumber}{xxx} %número da última página em algarismos romanos
\newcommand{\myTotalPagesNumber}{xxxxx} %quantidade total de páginas

\newcommand{\myRefereeONE}{Prof. Dr. xxxx}
\newcommand{\myRefereeTWO}{Profª. Drª. xxxxx}
\newcommand{\myRefereeTHREE}{Profª. Drª. xxxxx}
%Usar Profª. Drª. para o sexo feminino.

\newcommand{\deptChair}{Prof. Dr. xxxxx}  %nome do chefe de departamento
\newcommand{\programCoord}{Prof. Dr. xxxxxxxxxxxx} %nome do Coordenador do PROFMAT
%Usar Profª. Drª. para o sexo feminino.

%Definir até 5 palavras-chave em ordem de relevância. Não retirar o número e nem o ponto final. Apagar tudo que estiver entre chaves se não for utilizar as 5 palavras-chave, por exemplo {5. Palavra. \xspace} virará { }
\newcommand{\myKeyWordONE}{1. xxxxxx. \xspace}
\newcommand{\myKeyWordTWO}{2. xxxxxx. \xspace}
\newcommand{\myKeyWordTHREE}{3. xxxxxx. \xspace}
\newcommand{\myKeyWordFOUR}{4. xxxxxx. \xspace}
\newcommand{\myKeyWordFIVE}{5. xxxxxx. \xspace}

%traduza na mesma ordem as palavras-chave
\newcommand{\myKeyWordONEeng}{1.xxxxxx. \xspace}
\newcommand{\myKeyWordTWOeng}{2. xxxxxx.\xspace}
\newcommand{\myKeyWordTHREEeng}{3. xxxxxx.  \xspace}
\newcommand{\myKeyWordFOUReng}{4. xxxxxx.  \xspace}
\newcommand{\myKeyWordFIVEeng}{5.xxxxxx.  \xspace}
%%%%%%%%%%%%%%%%%%%%%%%%%%%%%%%%%%%%%%%%%%%%%%%%%%%%%%%%
%%%%%%%%%%%%%%%%%%%%%%%%%%%%%%%%%%%%%%%%%%%%%%%%%%%%%%%%
% A partir daqui o texto é gerado automaticamente 
% com os dados preenchidos acima.
%%%%%%%%%%%%%%%%%%%%%%%%%%%%%%%%%%%%%%%%%%%%%%%%%%%%%%%%
%%%%%%%%%%%%%%%%%%%%%%%%%%%%%%%%%%%%%%%%%%%%%%%%%%%%%%%%

	
\begin{document}
		
\pagenumbering{roman}
\thispagestyle{empty}

\hfill

\vspace{-2.0cm}

\begin{center}
	\begin{tabular}{c}
		\hline\\
		{\LARGE  \textbf{Universidade Federal de São Paulo}} \\ \\ 
		{\LARGE Instituto de Ciência e Tecnologia } \\ \\ 
		\hline
	\end{tabular}
\end{center}

\begin{center}
	\includegraphics[width=4cm]{Logotipo_UNIFESP.png} \\ \medskip %figure
\end{center}

\vfill

\begin{center}
	{\LARGE\textbf{Mestrado Profissional em  Matemática} }
\end{center}
\begin{center}
	{\LARGE \textbf{em Rede Nacional - PROFMAT}}
\end{center}
\vfill

\begin{center}
	{ \LARGE \textbf \myTitle }
\end{center}

\vfill


\begin{center}
	{\Large{\textbf{\myName}}}
\end{center}

\vspace{0.1cm}

\begin{center}
{\Large Orientador(a): \myProf}
\end{center}

\vfill

\begin{center}
	\textbf{\myLocation  \\  \myTime}
\end{center}

\newpage
\thispagestyle{empty}

\hfill

\vfill

\begin{center}
	\includegraphics[width=3.8cm]{Prof-Mat.png} \\ 
\end{center}


\noindent Título: \textit{\myTitle}

\noindent
Dissertação apresentada ao Instituto de Ciência e Tecnologia da UNIFESP, campus São José dos Campos/SP, como parte dos requisitos exigidos para a obtenção do título de Mestre pelo Programa de Mestrado Profissional em Matemática em Rede Nacional -- PROFMAT. 

\begin{center}
	\textbf{\myLocation  \\  \myTime}
\end{center}



\newpage
\thispagestyle{empty}

\quad

\vfill

\fbox{ \small \begin{minipage}[c]{13cm}
\begin{minipage}{12.5cm} % 12.5cm
			
			\
			
			\
			
			\
			
			\ \ \ \ \myNameSN
			
			\hspace{1cm} \textbf{\myTitle}, \myName \ -- São José dos Campos, \myConclusionYear.
			
			\hspace{1cm} \myLastRomanPageNumber, \myTotalPagesNumber f.
			
			\vspace{0.5cm}
			\hspace{1cm}  Dissertação (Mestrado) -- Universidade Federal de São Paulo. Instituto de Ciência e Tecnologia. Programa de Pós-Graduação em Matemática em Rede Nacional  (PROFMAT).
			
			\vspace{0.5cm}
			\hspace{1cm} \myEnglishTitle
			
			\vspace{0.5cm}
			\hspace{1cm} \myKeyWordONE \myKeyWordTWO \myKeyWordTHREE \myKeyWordFOUR \myKeyWordFIVE
			
			\
			
			\
			
		\end{minipage}
		\end{minipage}
} 

		
	
	

	\newpage
	\thispagestyle{plain}
	
	
	\begin{center}
		\textbf{UNIVERSIDADE FEDERAL DE SÃO PAULO }
		
		\vspace{.3cm}
		
		\textbf{INSTITUTO DE CIÊNCIA E TECNOLOGIA} 
		
		\vspace{.6cm}    	   
		
		\textbf{Mestrado Profissional em Matemática em Rede Nacional}
		
		\textbf{PROFMAT}
	\end{center}
	
	\vfill 
	
	\textbf{Chefe de departamento:}     
	
	\hspace{1cm} \deptChair 
	
	\vspace{.3cm} 
	
	\textbf{Coordenadora do Programa de Pós-Graduação:} 
	
	\hspace{1cm} \programCoord 
	
	\vfill                      
	\newpage
	\thispagestyle{plain}
	
	\vfill
	
	\begin{center}
		\myName
		
		\vspace{.3cm}
		
		\myTitle 
		
	\end{center}
	
	
	
	\vfill
	
	
	\textbf{Presidente da banca:} \myProf \medskip
	
	\vspace{.3cm} 
	
	\textbf{Banca examinadora:} \medskip 
	
	\myRefereeONE \medskip
	
	\myRefereeTWO \medskip
	
	\myRefereeTHREE \medskip
	
	\vspace{.3cm} 
	
	\textbf{Data da Defesa:} \myDefenseDay
	
	\vfill
	
	
	
	%%%%%%%%%%%%%%%%%%%%%%%%%%%%%%%%%%%%%%%%%%%%%%%%%%%%%%%%
	%%%%%%%%%%%%%%%%%%%%%%%%%%%%%%%%%%%%%%%%%%%%%%%%%%%%%%%%
	% A partir daqui o aluno passa a fazer o trabalho\ldots
	%%%%%%%%%%%%%%%%%%%%%%%%%%%%%%%%%%%%%%%%%%%%%%%%%%%%%%%%
	%%%%%%%%%%%%%%%%%%%%%%%%%%%%%%%%%%%%%%%%%%%%%%%%%%%%%%%%
	
	\newpage
	\thispagestyle{plain}

\hfill

\vfill

\begin{flushright}
	{\slshape
	``Não é apenas a questão, mas a maneira como você tenta resolvê-la.'' 
	\\ Maryam Mirzakhani
	} %inserir uma frase. Este é apenas um exemplo
\end{flushright}


\chapter*{Agradecimentos}






\refstepcounter{dummy}
\newpage
\begingroup
\let\clearpage\relax
\let\cleardoublepage\relax
\let\cleardoublepage\relax
\chapter*{Resumo}
\noindent








\

\noindent
\textbf{Palavras-chave:} \myKeyWordONE \myKeyWordTWO \myKeyWordTHREE \myKeyWordFOUR \myKeyWordFIVE 

	
	
\newpage



\chapter*{Abstract}
\noindent




\noindent
\textbf{Keywords:} \myKeyWordONEeng \myKeyWordTWOeng \myKeyWordTHREEeng \myKeyWordFOUReng \myKeyWordFIVEeng
\endgroup			

\vfill


\mainmatter
\tableofcontents 


\chapter*{Introdução}
\addcontentsline{toc}{chapter}{Introdução}

\chapter{Título}

\chapter{Título}    
  
    


  





\begin{thebibliography}{56}

%\bibitem{barra1} Automaclick. Como Saber de Onde Vem o Produto Pelo Código de Barras?. (2017)

%\url{https://www.automaclick.com.br/blog/saber-de-onde-vem-o-produto-pelo-codigo-de-barras}. Acesso em: 02/09/2023

\bibitem{bncc} BRASIL. Ministério da Educação. Base Nacional Comum Curricular. Brasília,  (2017).

\url{http://basenacionalcomum.mec.gov.br/images/BNCC_EI_EF_110518_versaofinal_site.pdf}.

\bibitem{coutinho}  Coutinho, Severino Collier. Criptografia. Programa de Iniciação Científica da OBMEP. IMPA. (2015).


%\bibitem{caldeira} Caldeira, Johne; Matos, Rayane Auxiliadora de Oliveira. Códigos Corretores de Erros e uma Breve Introdução aos Códigos de Hamming. Revista Focco, (2023).

%\url{file:///C:/Users/ACER-A515-Ssd/Downloads/01+FOCO+-+DOI%20(4).pdf} Acesso em: 21/09/2023.

%\bibitem{barra2} Códigos de Barras Brasil. A Evolução e a História do Código de Barras.

%\url{https://codigosdebarrasbrasil.com.br/a-evolucao-e-historia-do-codigo-de-barras/}. Acesso em 19/08

%\bibitem{barra3} Códigos de Barras Brasil. Códigos de Barras: O que é, Como Funcionam e Principais tipos.

%\url{https://codigosdebarrasbrasil.com.br/tipos-de-codigos-de-barras/}. Acesso em 13/08

\bibitem{lais} de Melo, Lais Aline Casagrande Pires. Decifrando a Aritmética para o Ensino Fundamental. Dissertação de Mestrado. Mestrado Profissional em Matemática em Rede Nacional - PROFMAT. UNIFESP. São José dos Campos. (2017).

\url{https://sca.profmat-sbm.org.br/profmat_tcc.php?id1=3196&id2=150870749}. Acesso em 02/12/23.

\bibitem{george} de Morais, George Absalão Pandino. Códigos Perfeitos na Métrica de Lee e a Conjectura de Golomb-Welch. Mestrado em Matemática Aplicada. UNIFESP. São José dos Campos. (2017).

\url{https://repositorio.unifesp.br/items/a1c2d791-8c84-4633-be40-2a6e0047dd2d}. Acesso em 27/05/24.


%\bibitem{castro} Castro, Jânio Kléo Sousa. Teoria dos Números / Jânio Kléo Sousa de Castro; Coordenação Cassandra Ribeiro Joye - Fortaleza: UAB/IFCE (2010).
%\url{https://www.educapes.capes.gov.br/bitstream/capes/429385/2/Teoria%20dos%20Numeros.pdf}

\bibitem{faleiros} Faleiros, Antonio Cândido. Criptografia. Notas em Matemática Aplicada. SBMAC.  (2011).

\url{https://www.sbmac.org.br/wp-content/uploads/2022/08/livro_52.pdf}. Acesso em 27/05/24.

\bibitem{fiarresga} Fiarresga, Victor Manuel Calhabrês. Criptografia e matemática.  Lisboa, Portugal. 161 p. Tese (Mestrado) - Faculdade de Ciências, Universidade de Lisboa  (2010). 
\url{https://repositorio.ul.pt/bitstream/10451/3647/1/ulfc055857_tm_Victor_Fiarresga.pdf} Acesso em: 15/05/2024.

\bibitem{abramoprofmat} Hefez, Abramo. Aritmética. SBM. Coleção PROFMAT. (2014).

\bibitem{abramo1} Hefez, Abramo. Iniciação à Aritmética. Rio de Janeiro: IMPA, (2015).

\url{http://www.obmep.org.br/docs/apostila1.pdf}.

\bibitem{abramo2} Hefez, Abramo; Villela, T. Maria Lúcia. Códigos Corretores de Erros. Rio de Janeiro: IMPA, (2008).

%\bibitem{isbn} ISBN Intrenational Agency. O que é um ISBN?.(2014)

%\url{https://www.isbn-international.org/content/what-isbn/10}. Acesso em 13/09/2023.

\bibitem{iusenko} Iusenko, Konstiantyn. Álgebra I para Licenciatura, Notas de aula, IME, USP, São Paulo, 2021.

\url{https://www.ime.usp.br/~iusenko/ensino_2021_1/index.php}. Acesso em 08/09/23.

\bibitem{jordão} Jordão, Edmarcos Martins. RSA-Crypta: Uma Aplicação da Criptografia no Ensino Médio. Dissertação de Mestrado. Mestrado Profissional em Matemática em Rede Nacional - PROFMAT. Universidade Federal Vale do São Francisco. Juazeiro. (2023).

\url{https://sca.profmat-sbm.org.br/profmat_tcc.php?id1=7171&id2=171055450}.

\bibitem{sueli} Lavor, Carlile Campos; Alves, Marcelo Muniz Silva; de Siqueira, Rogério Monteiro; Costa, Sueli Irene Rodrigues. Uma Introdução à Teoria de Códigos. Notas em Matemática Aplicada. SBMAC. (2012).

\url{https://www.sbmac.org.br/wp-content/uploads/2022/08/livro_21.pdf}.

\bibitem{manoel} Lemos, Manoel. Criptografia, Números Primos e Algoritmos. Rio de Janeiro, IMPA. (2010).
\url{https://impa.br/wp-content/uploads/2017/04/PM_04.pdf}. Acesso em 17/05/24.

\bibitem{morgado} Morgado, Augusto Cezar de Oliveira; Carvalho, Paulo Cezar Pinto Carvalho. Matemática Discreta. SBM. Coleção PROFMAT. (2015).


\bibitem{nicoletti} Nicoletti, Everton Rodrigo. Aplicações de Álgebra Linear aos Códigos Corretores de Erros e ao Ensino Médio. Mestrado Profissional em Matemática em Rede Nacional - PROFMAT, UNESP, Rio Claro. (2015).

\url{https://sca.profmat-sbm.org.br/profmat_tcc.php?id1=1638&id2=507}. Acesso em: 13/04/2023.

\bibitem{roberta} Ribeiro, Roberta Alves do Nascimento. Os Códigos de Barras como Temas Geradores no Ensino de Matemática. UNESP. Faculdade de Engenharia e Ciências de Guaratinguetá, (2021). 

\url{https://repositorio.unesp.br/server/api/core/bitstreams/89d6a986-0dbc-4daa-a275-8afb5fbee3c6/content}. Acesso em: 05/05/2023.

\bibitem{graça} Sabadin, Graça Aparecida Prestes. Códigos Corretores de Erros. Mestrado Profissional em Matemática em Rede Nacional. UFSC. (2019).

\url{https://sca.profmat-sbm.org.br/profmat_tcc.php?id1=4936&id2=170660171}. Acesso em: 13/04/2023.

\bibitem{paulo} Santos, Paulo César. Códigos Verificadores e  Corretores de Erros. Mestrado Profissional em Matemática em Rede Nacional. UFABC. Santo André. (2018).

\url{https://sca.profmat-sbm.org.br/profmat_tcc.php?id1=4256&id2=150530085}. Acesso em: 09/05/2023.


 \url{https://cdnportaldaobmep.impa.br/portaldaobmep/uploads/material_teorico/83bhrw1mjmgwo.pdf}. Acesso em: 30/11/23.
 
\bibitem{winterle} Steinbruch, Alfredo; Winterle, Paulo. Álgebra Linear 2. São Paulo: Pearson/Makron Books. (1987).

\bibitem{wikipedia1} Wikipédia, A Enciclopédia Livre. Mariner 9.(2022)

\url{https://gl.wikipedia.org/wiki/Mariner_9}. Acesso em 07/09/2023.

\bibitem{wikipedia2} Wikipédia, A Enciclopédia Livre. Voyager 1.(2022).

\url{https://pt.wikipedia.org/wiki/Voyager_1}. Acesso em 07/09/2023.


\end{thebibliography}







              \end{document}

\end{document}